% ===================================================================
%  BẢNG Số 10 — Biển. Cảng.
%
%  Traduction, faite en 2026, en vietnamien standard d'aujourd'hui,
%  sur l'ido de texto/io. Regles de la colonne : voir 10-bang-01.tex.
%
%  LE TABLEAU DE LA MARINE EST CELUI OU LE VIETNAMIEN A LE PLUS DE
%  MOTS A LUI, ET C'EST L'INVERSE DE CE QU'ON ATTENDAIT. Au Quebec,
%  ce meme tableau etait le moins regional de tous, parce que le
%  vocabulaire de la marine est un vocabulaire de metier fixe par la
%  meme ecole des deux cotes de l'Atlantique. Ici il est le plus
%  vietnamien, et pour la raison exactement contraire : un pays de
%  trois mille kilometres de cote avait ses mots avant d'avoir ses
%  manuels. Thuyền la barque, buồm la voile, chèo la rame, lưới le
%  filet, neo l'ancre, bến le quai, cảng le port, sóng la vague,
%  thủy triều la maree, hải đăng le phare, bãi biển la plage — pas un
%  emprunt.
%
%  ET CE QUI EST EMPRUNTE EST CE QUI EST ARRIVE PAR L'ARSENAL : « tàu
%  ngầm » le sous-marin, « tàu tuần dương » le croiseur, « ụ tàu » le
%  bassin de radoub, « ca-nô » le canot a moteur, « xà lan » le
%  chaland. Le militaire et l'industriel sont sino-vietnamiens ou
%  francais ; le pecheur et le batelier sont vietnamiens.
%
%  UNE SEULE HESITATION, ET ELLE VAUT D'ETRE DITE : l'ido ecrit
%  « warfo », qui est l'anglais \textit{wharf}, la ou il ecrit
%  « kayo » pour le quai. Le quebecois avait garde « jetee » parce
%  que « quai » etait deja pris ; le vietnamien fait de meme avec
%  « đê chắn sóng », la digue qui arrete les vagues, et « bến » pour
%  le quai. Les deux ouvrages restent distincts et portent chacun son
%  renvoi.
% ===================================================================

%%K t10-apar-1 apar
\VUsaut{4.0mm}
\VUtitre{15.0pt}{60}{BẢNG Số 10}
\VUsaut{3.0mm}
\VUpk{11.4pt}{40}{Biển. --- Cảng.}
\VUsaut{3.0mm}

%%K t10-01-1 p
1. --- Mùa hè, trong khi người này tìm sự yên tĩnh và mát mẻ trên núi, người kia lại
thích ra bờ biển. Năm ngoái chúng em đã làm như thế, vì chúng em rất yêu \VUgras{Đại
dương} \textsuperscript{(1)}.

%%K t10-02-1 p
2. --- Riêng em, em thấy một niềm vui lớn khi ngắm khoảng nước mênh mông trải đến tận
\VUgras{chân trời} \textsuperscript{(2)}, và khi nhìn những chiếc \VUgras{tàu hơi nước}
\textsuperscript{(3)} khổng lồ chở khách đi \VUgras{vượt biển} sang châu Mỹ.

%%K t10-03-1 p
3. --- Vì thế cha em, người biết sở thích của em, bao giờ cũng chọn cho kỳ nghỉ một bãi
biển thật yên tĩnh, gần một trong những \VUgras{quân cảng} lớn của nước ta.

%%K t10-03-2 p
Trong kỳ nghỉ vừa rồi, \VUgras{hạm đội} đến thả neo sau \VUgras{đê chắn sóng}
\textsuperscript{(4)} khổng lồ đóng kín và che chở \VUgras{quân cảng}
\textsuperscript{(5)}, và em đã được thỏa lòng ngắm các loại \VUgras{tàu chiến} khác
nhau, từ chiếc \VUgras{tàu ngầm} \textsuperscript{(10)} nhỏ lặn xuống và biến mất dưới
nước, đến chiếc \VUgras{thiết giáp hạm} \textsuperscript{(9)} khổng lồ, thứ mà đến nay
hầu như chỉ sợ những đòn tấn công của \VUgras{tàu phóng lôi} \textsuperscript{(8)}.

%%K t10-tit-2 sub
\VUsaut{3.0mm}
\VUpk{10.2pt}{40}{Những con tàu nhỏ.}
\VUsaut{3.0mm}

%%K t10-04-1 p
4. --- Chiếc tàu phóng lôi ấy đã biết tìm rất \VUgras{khéo} điểm yếu của lớp giáp kim
loại dày để cắm vào đó quả \VUgras{ngư lôi} nguy hiểm, mà tiếng nổ làm con tàu chìm;
nhưng dù sao nó vẫn ít đáng sợ hơn tàu ngầm, vì tàu săn ngư lôi dễ dàng đuổi theo nó.

%%K t10-05-1 p
5. --- Em cũng được thấy những chiếc \VUgras{tàu tuần dương} \textsuperscript{(6)} bọc
thép nhanh nhẹn, gần như bất khả xâm phạm chẳng kém thiết giáp hạm thật, mấy chiếc
\VUgras{tàu vận tải} \textsuperscript{(12)}, ba bốn chiếc \VUgras{pháo hạm}
\textsuperscript{(7)} và sau cùng một chiếc \VUgras{khinh hạm} \textsuperscript{(11)}.
\VUgras{Cảnh hạm đội ấy nhổ neo là cảnh thú vị nhất.}

%%K t10-06-1 p
6. --- Cách đóng khác nhau biết bao giữa những khối thép lớn ấy với những chiếc
\VUgras{tàu ba cột buồm} thanh nhã hay những chiếc \VUgras{thuyền buồm}
\textsuperscript{(13)} mảnh mai mà đội thương thuyền còn dùng đến ngày nay, và em đã thấy
chúng vào cảng, giương hết buồm, trong lúc em dạo trên \VUgras{đê}
\textsuperscript{(14)}. Thật ra chúng mất nhiều lợi thế \VUgras{khi} \VUgras{gió} ngược.
Chúng phải hạ hết buồm để vào lại bến, và cần một chiếc \VUgras{tàu kéo}
\textsuperscript{(16)} ra dắt chúng vào bến như những chiếc \VUgras{xà lan}
\textsuperscript{(17)} thường.

%%K t10-tit-3 sub
\VUsaut{3.0mm}
\VUpk{10.2pt}{40}{Thương cảng.}
\VUsaut{3.0mm}

%%K t10-07-1 p
7. --- Điều thú vị ở cái cảng em đang nói đến là nó không chỉ có một quân cảng, được xây
dựng nhân tạo bằng những công trình khổng lồ, mà còn có một \VUgras{thương cảng} tuyệt
vời nằm ở \VUgras{cửa} một con sông lớn.

%%K t10-08-1 p
8. --- Dải đất ngăn hai vùng bến được xây thành công sự và bảo vệ bằng một loạt
\VUgras{ụ pháo} và \VUgras{pháo đài nhỏ} nhô ra biển. Muốn dạo trên \VUgras{tường thành}
\textsuperscript{(15)} phải có giấy phép riêng. Em đã xin được dễ dàng nhờ bác em, kỹ sư
của \VUgras{xưởng đóng tàu} \textsuperscript{(18)}, và em đã đi xem những con tàu đang
được đóng dưới các nhà xưởng rộng lớn.

%%K t10-09-1 p
9. --- Em còn được dự lễ hạ thủy một chiếc thiết giáp hạm, và em nhớ giây phút xúc động
mà cả đám đông cảm thấy khi khối sắt khổng lồ, được thả ra, bắt đầu trượt trên bệ đỡ và
nổi lên mặt nước sông.

%%K t10-10-1 p
10. --- Muốn đến xưởng đóng tàu, người ta đi qua \VUgras{bãi tập}
\textsuperscript{(19)}, nơi lính đồn trú tập luyện mỗi ngày, rồi qua một \VUgras{cầu
nhỏ} \textsuperscript{(20)} bằng sắt nối sân duyệt binh với \VUgras{xưởng vũ khí}
\textsuperscript{(21)}.

%%K t10-tit-4 sub
\VUsaut{3.0mm}
\VUpk{10.2pt}{40}{Xưởng vũ khí và các ụ tàu.}
\VUsaut{3.0mm}

%%K t10-11-1 p
11. --- Cùng bác em, em đã thăm cơ sở lớn ấy, đầy \VUgras{đại bác}
\textsuperscript{(22)}, \VUgras{đạn tròn} \textsuperscript{(23)} và \VUgras{đạn trái
phá}. Em cũng thấy \VUgras{xưởng luyện kim} \textsuperscript{(24)}, nơi người ta chế tạo
\VUgras{nồi hơi} \textsuperscript{(25)} cho tàu hơi nước.

%%K t10-12-1 p
12. --- Ngay gần đó là những \VUgras{ụ sửa tàu} \textsuperscript{(26)} và những
\VUgras{ụ nổi} \textsuperscript{(27)} rất lạ mắt, nơi em đã được thấy thật gần, trên một
con tàu đang sửa, chiếc \VUgras{chân vịt} \textsuperscript{(28)}, cái \VUgras{sống tàu}
\textsuperscript{(29)} và cái \VUgras{bánh lái} \textsuperscript{(30)}.

%%K t10-13-1 p
13. --- Thương cảng cũng đáng nghiên cứu chẳng kém quân cảng, và em đã đứng hàng giờ
ngắm nghía trên các bến, nơi buộc những chiếc \VUgras{tàu bánh guồng}
\textsuperscript{(31)} ngược dòng sông, và nhìn hoạt động của những chiếc \VUgras{cần
cẩu dựng cột buồm} \textsuperscript{(32)} nhấc bổng như lông hồng những kiện hàng khổng
lồ.

%%K t10-tit-5 sub
\VUsaut{4.0mm}
\VUpk{10.2pt}{40}{Hoa tiêu.}
\VUsaut{3.0mm}

%%K t10-14-1 p
14. --- Em thán phục những chiếc \VUgras{tàu vượt Đại Tây Dương}
\textsuperscript{(34)} ra khỏi vũng đậu, \VUgras{cờ hiệu} \textsuperscript{(35)} tung
bay, do một \VUgras{hoa tiêu} \textsuperscript{(36)} tài giỏi dẫn đường, người biết đi
theo \VUgras{luồng} một cách kỳ diệu nhờ những chiếc \VUgras{phao}
\textsuperscript{(40)} và \VUgras{tiêu} đánh dấu, tránh những chiếc \VUgras{thuyền chở
hàng} \textsuperscript{(41)} nhọc nhằn lái đi bằng cánh buồm vuông duy nhất của chúng. Em
phân biệt được ở \VUgras{phía trước} \textsuperscript{(37)} --- mũi tàu --- những
\VUgras{ca-bin} \textsuperscript{(39)} hạng hai, và ở \VUgras{phía sau}
\textsuperscript{(38)} --- đuôi tàu --- những phòng khách của hành khách hạng nhất.

%%K t10-15-1 p
15. --- Đôi khi em cũng dự việc chất hàng lên một chiếc \VUgras{tàu chở hàng}
\textsuperscript{(42)}, nơi người ta vừa chất hàng hóa từ \VUgras{kho}
\textsuperscript{(43)} và từ \VUgras{ga cảng} \textsuperscript{(44)}, do vô số chuyến tàu
của \VUgras{tuyến đường sắt bến cảng} \textsuperscript{(45)} chở đến.

%%K t10-tit-6 sub
\VUsaut{3.0mm}
\VUpk{10.2pt}{40}{Chèo thuyền.}
\VUsaut{3.0mm}

%%K t10-16-1 p
16. --- Em thường chèo thuyền trong \VUgras{bến nổi} \textsuperscript{(53)}, nơi những
chiếc \VUgras{thuyền nhỏ} \textsuperscript{(46)} vào trú, và em vui sướng cầm cái
\VUgras{mái chèo} \textsuperscript{(47)}. Em leo lên \VUgras{boong}
\textsuperscript{(48)} một chiếc \VUgras{thuyền buồm} \textsuperscript{(49)}, em bám
\VUgras{dây chão} \textsuperscript{(52)} trèo dọc \VUgras{cột buồm}
\textsuperscript{(51)} lên tận những \VUgras{trục buồm} \textsuperscript{(50)}, rồi em
lại đi dạo bằng \VUgras{thuyền đua} \textsuperscript{(54)} giữa những chiếc \VUgras{thuyền
đánh cá} \textsuperscript{(55)} nặng nề, nơi \VUgras{lưới} \textsuperscript{(56)} đang
phơi khô.

%%K t10-17-1 p
17. --- Trên \VUgras{bến} \textsuperscript{(58)}, trước mắt em diễu qua đủ loại
\VUgras{hàng hóa} \textsuperscript{(59)}, đựng trong những \VUgras{thùng}
\textsuperscript{(60)} hay những \VUgras{kiện} \textsuperscript{(61)}. Chúng làm thành
\VUgras{hàng chở} \textsuperscript{(63)} của các \VUgras{sà lúp}, và những chiếc
\VUgras{cần cẩu} \textsuperscript{(62)} khỏe nhấc chúng lên và đặt xuống những chiếc
\VUgras{xe tải} \textsuperscript{(64)}, trong khi \VUgras{người chở hàng} khai báo và trả
thuế ở \VUgras{sở quan thuế} \textsuperscript{(65)}.

%%K t10-18-1 p
18. --- Sau cùng, khi đến \VUgras{cầu quay} \textsuperscript{(66)} hoặc \VUgras{âu tàu}
\textsuperscript{(69)}, đi ngang chiếc \VUgras{tàu cuốc} \textsuperscript{(68)} đang nạo
\VUgras{con kênh} \textsuperscript{(67)}, em lên bờ ở đê, cạnh \VUgras{cột tín hiệu}
\textsuperscript{(70)}.

%%K t10-19-1 p
19. --- Chính ở phía bên kia của con đê ấy mà những chiếc \VUgras{xuồng}
\textsuperscript{(71)} chở sĩ quan của hạm đội cập vào bờ. Thật đáng ngắm cảnh các thủy
thủ nhịp nhàng nâng mái chèo, trong khi con thuyền, được \VUgras{sóng}
\textsuperscript{(72)} nâng lên, dễ dàng cập vào bậc thang lên bờ. Chính ở đó người ta
quan sát rõ nhất \VUgras{thủy triều} cùng sự lên xuống của \VUgras{nước lớn} và
\VUgras{nước ròng} trên \VUgras{bãi biển} \textsuperscript{(73)}.

%%K t10-tit-7 sub
\VUsaut{3.0mm}
\VUpk{10.2pt}{40}{Câu lạc bộ.}
\VUsaut{3.0mm}

%%K t10-20-1 p
20. --- Để về nhà, em đi \VUgras{xe điện} \textsuperscript{(74)}, có \VUgras{trạm dừng}
\textsuperscript{(75)} gần cái \VUgras{cột} dựng trước câu lạc bộ sĩ quan.

%%K t10-20-2 p
Từ \VUgras{ban công} \textsuperscript{(76)} của câu lạc bộ, trang trí bằng những chiếc
\VUgras{mỏ neo} \textsuperscript{(77)}, đại bác và đạn tròn, em thường thấy một cuộc tụ
họp lộng lẫy của các sĩ quan hải quân : những \VUgras{đại úy hải quân}
\textsuperscript{(78)} đeo kiếm, những \VUgras{đại úy pháo binh}
\textsuperscript{(79)} và \VUgras{bộ binh} \textsuperscript{(80)} đeo \VUgras{gươm}.
Đứng sau họ là một \VUgras{thủy thủ} \textsuperscript{(81)} mang \VUgras{ống nhòm} đến
cho họ khi họ muốn nhìn rõ những gì xảy ra ở xa.

%%K t10-21-1 p
21. --- Em cũng thấy những \VUgras{thiếu úy} \textsuperscript{(82)} trẻ nhận lệnh của
\VUgras{hạm trưởng} \textsuperscript{(83)} hay của \VUgras{đô đốc}
\textsuperscript{(84)}, trong khi một \VUgras{hạ sĩ quan} \textsuperscript{(85)} chờ để
mang lệnh ra tàu.

%%K t10-22-1 p
22. --- Từ ban công ấy, cảnh nhìn thật tuyệt : người ta bao quát cả \VUgras{bãi biển}
với những chiếc \VUgras{lều} \textsuperscript{(86)} và \VUgras{buồng thay đồ}
\textsuperscript{(87)}, và vào giờ tắm, người ta thấy những \VUgras{người tắm biển}
\textsuperscript{(88)} nô đùa vui vẻ trước \VUgras{sòng bạc} \textsuperscript{(89)}.

%%K t10-23-1 p
23. --- Ở cuối bãi biển, bờ biển tạo thành một \VUgras{bán đảo}
\textsuperscript{(91)} nhỏ đầy \VUgras{đá}, nơi \VUgras{sóng lớn} \textsuperscript{(92)}
vỗ vào và trên đó dựng một ngọn \VUgras{hải đăng} \textsuperscript{(93)} chỉ đường ban
đêm cho tàu bè. Bán đảo ấy chỉ nối với đất liền bằng một \VUgras{eo đất} rất hẹp
\textsuperscript{(90)}.

%%K t10-tit-8 sub
\VUsaut{3.0mm}
\VUpk{10.2pt}{40}{Toàn cảnh bờ biển.}
\VUsaut{3.0mm}

%%K t10-24-1 p
24. --- \VUgras{Vách đá} \textsuperscript{(94)} trải dài, bị biển khoét thành một loạt
\VUgras{vịnh} \textsuperscript{(95)} và \VUgras{mũi đất}, làm cho nó thật nên thơ.

%%K t10-24-2 p
Trên \VUgras{cao nguyên} \textsuperscript{(96)} nhìn xuống \VUgras{eo biển}
\textsuperscript{(98)} có một \VUgras{pháo đài} \textsuperscript{(97)} rất quan trọng và
mấy tòa « biệt thự ».

%%K t10-25-1 p
25. --- Eo biển ấy ngăn với đất liền một \VUgras{hòn đảo} \textsuperscript{(99)} rất
nguy hiểm, chung quanh đầy \VUgras{đá ngầm} \textsuperscript{(100)} và \VUgras{bãi sóng
vỡ}, nơi thuyền bè và cả những con tàu lớn đôi khi mắc cạn. Dù cứu hộ có nhanh và
\VUgras{xuồng cứu nạn} đến kịp, những vụ \VUgras{đắm tàu} ấy vẫn thật khủng khiếp, và
trong những \VUgras{cơn bão} mùa phân, nhiều con tàu đã chìm cùng người và hàng.

%%K t10-25-2 p
Dù có mối nguy bên cạnh, \VUgras{cảng} vẫn rất đông thủy thủ lui tới.
\VUsaut{6.0mm}
\VUfilet{22mm}
